% ==========================================================
% == Section 6: Conclusion                                 ==
% ==========================================================
\section{Conclusion}
\label{sec:conclusion}

% [段 1 · 核心贡献重述] 大意:回扣 intro 的问题框架,总结 DACE 把"覆盖广度"和"记忆持久度"作为协同演化的两个新维度。呼应 abstract 和 intro 最后的贡献 bullet。
% 撰写逻辑:先复述问题—"协同演化只保证双方进步,但不保证广度和持久度";再概括 DACE 的两个核心机制—归一化覆盖增益与 Beta-Bernoulli 贝叶斯回放;最后用一句话概括实验结论:防御 ASR 更低、通用能力不降、攻击覆盖更广、并出现组合攻击涌现。
% 中文对应内容:我们提出 DACE,一个以多样性驱动的协同演化框架,旨在解决 RL 协同安全训练中的两大病态——攻击策略坍缩与防御对抗遗忘。DACE 的两大机制分别对应这两个病态:显式的 12×10 攻击策略空间配合归一化边际覆盖增益奖励保证攻击的结构化广度与信号稳定性;Beta-Bernoulli 贝叶斯对抗回放池配合时间衰减与 Thompson 采样在连续不确定性度量下动态维持防御记忆。主实验证明 DACE 在不牺牲通用能力的前提下降低了多个 benchmark 上的 ASR,显著提高了攻击策略覆盖率,并涌现出组合式攻击这一新现象。
This paper addresses a fundamental limitation of reinforcement-learning-based co-evolutionary safety alignment: standard frameworks guarantee that both the attacker and the defender \emph{keep improving}, but not that attacks \emph{cover} the true threat space or that defenses \emph{remain} robust to past threats. We name and formalize these two compounding pathologies---attack strategy collapse and defense adversarial forgetting---and propose \method, which addresses them jointly with two principled mechanisms: (i) an explicit $12\times10$ risk-by-style attack strategy grid combined with a \emph{normalized marginal coverage gain} reward that provides a bounded, non-vanishing signal toward under-explored strategies and admits a closed-form KL-contraction interpretation; and (ii) a unified \emph{Bayesian adversarial replay pool} in which every successful attack carries a Beta--Bernoulli threat posterior subject to time decay and Thompson sampling, so that defense training is continuously anchored to the defender's evolving threat landscape rather than to a stale discrete tally. Empirically, \method improves safety on both in-distribution and OOD-attacker benchmarks, preserves general capabilities, substantially broadens attack strategy coverage, and produces a new qualitative phenomenon---combinatorial attack emergence---that is absent in prior co-evolutionary baselines.

% [段 2 · 局限] 大意:诚实列出三类局限,为未来工作铺路,也让 review 知道我们已意识到边界条件。
% 撰写逻辑:三类局限串在一段内,分别对应 (a) 策略空间粒度局限、(b) 依赖 SFT 质量与判官稳定性、(c) 本文仅验证了单轮文本设置。
% 中文对应内容:DACE 仍有若干局限。首先,12×10 的策略空间是离散且由人工设计的,既可能遗漏跨类别的风险,也可能被未来攻击者学会绕过编号;未来可以探索自动或连续的策略空间。其次,攻击者的组合涌现对 SFT 热启动质量与 judge 模型的稳定性敏感,若上游数据或 judge 存在偏置,这些偏置可能被协同演化放大。第三,本文的实验仅覆盖了文本单轮场景,DACE 向多模态、工具使用与多轮 agent 的扩展尚待系统验证。
\paragraph{Limitations.} Several limitations of \method remain. \emph{First}, the $12\times10$ strategy space is discrete and manually curated; while our ablations indicate the pruned grid is well-justified for a text-only defender, future attackers may exploit behaviors that straddle categories or lie outside the grid, and automated or continuous strategy spaces are a natural next step. \emph{Second}, combinatorial emergence and the quality of Beta--Bernoulli posteriors depend on the stability of the safety judge and the warm-start SFT data; upstream biases in either component may be amplified by co-evolution if not monitored. \emph{Third}, our evaluation is confined to text-only, single-turn interaction; extending \method to multimodal, tool-using, and multi-turn agentic settings is left to future work. A detailed discussion of failure modes and mitigations appears in Appendix~\ref{apdx:limitations}.

% [段 3 · 未来方向] 大意:给出方向性展望,呼应 intro 的"对抗持续演化的威胁需要同样持续演化的防御"主线,提示这一视角可以进一步推广。
% 撰写逻辑:三条方向,分别对应 (a) 扩展策略空间 / 自适应发现、(b) lifelong 对抗记忆跨 run 保存、(c) 与 inference-time 推理安全对齐的正交结合。最后一句上升到范式层面。
% 中文对应内容:未来工作可以从三个方向推进。其一,跨 RL 训练的 lifelong 对抗记忆,让档案池跨模型版本累积、并作为社区共享的策略-风险档案;其二,与 inference-time 推理安全对齐方法（如自省式 refusal、推理链审计)正交组合,覆盖 pre-train / post-train / test-time 全链条;其三,将"覆盖广度 + 记忆持久度"作为协同演化的通用设计原则,推广到工具使用、多 agent 协作、多模态安全对齐等更复杂的场景。我们相信把安全对齐视为"持续演化"而非"一次到位"将是下一阶段对齐研究的主线。
\paragraph{Future work.} Several directions follow naturally from \method. First, \emph{lifelong adversarial memory} that aggregates Beta--Bernoulli archives across RL runs and model generations could turn the replay pool into a community-shared, strategy-indexed threat registry. Second, \emph{orthogonal composition} of \method with inference-time safety alignment---e.g., introspective refusal reasoning~\citep{stair2025} or reasoning-chain auditing---could cover the pre-train / post-train / test-time pipeline end-to-end. Third, the broader principle of ``coverage $\times$ memory'' as co-evolutionary design objectives can be generalized to tool-using, multi-agent, and multimodal safety alignment, where the strategy space is typically richer and drift is typically faster. Taken together, these directions support a shift from ``one-shot'' alignment to \emph{alignment as continuous evolution}---and \method offers a concrete mechanism on which further such frameworks can build.
